\documentclass[12pt,a4paper]{report}
\usepackage[utf8]{inputenc}
\usepackage{amsmath}
\usepackage{amsfonts}
\usepackage{amssymb}
\usepackage{amsthm}
\usepackage{hyperref}

\usepackage{multicol}
\usepackage{fancyhdr}
\usepackage[inline]{enumitem}
\usepackage{tikz}
\usepackage{tikz-cd}
\usetikzlibrary{calc}
\usetikzlibrary{shapes.geometric}
\usepackage[margin=0.5in]{geometry}
\usepackage{xcolor}
\usepackage{ esint }

\hypersetup{
    colorlinks=true,
    linkcolor=blue,
    filecolor=magenta,      
    urlcolor=cyan,
    pdftitle={Tensors},
    pdfpagemode=FullScreen,
    }

%\urlstyle{same}

\newcommand{\CLASSNAME}{Math 5230 -- Partial Differential Equations}
\newcommand{\STUDENTNAME}{Paul Carmody}
\newcommand{\ASSIGNMENT}{NOTES }
\newcommand{\DUEDATE}{December 17, 2025}
\newcommand{\SEMESTER}{Fall 2025}
\newcommand{\SCHEDULE}{MW 2:00 -- 3:15}
\newcommand{\ROOM}{Crawford 332}

\newcommand{\MMN}{M_{m\times n}}
\newcommand{\FF}{\mathcal{F}}

\pagestyle{fancy}
\fancyhf{}
\chead{ \fancyplain{}{\CLASSNAME} }
%\chead{ \fancyplain{}{\STUDENTNAME} }
\rhead{\thepage}
\fancyfoot[C]{-- \thepage --}
\input{../MyMacros}

\newcommand{\RED}[1]{\textcolor{red}{#1}}
\newcommand{\BLUE}[1]{\textcolor{blue}{#1}}
\newcommand{\FFF}{\mathcal{F}}

\begin{document}

\begin{center}
	\Large{\CLASSNAME -- \SEMESTER} \\
	\large{ w/Professor Snelson}
\end{center}
\begin{center}
	\STUDENTNAME \\
	\ASSIGNMENT -- \DUEDATE\\
\end{center} 

\textbf{\Large{Basic PDEs}}
\begin{align*}
	u_x  &= 0 & \quad &\implies u(x,y) = f(y) \\
	u_{xx} &= 0 & \quad &\implies u(x,y) = f(y)x+g(y) \\
	u_{xx}-u &=0 & \quad &\implies u(x,y) = f(y)\cos x + g(y)\sin x \\
	u_{xy} &= 0 & \quad &\implies u(x,y) = F(y) + G(x)
\end{align*}for some functions $f(y)$ and $g(y)$.  Note: it is interesting that $u_{xy}$ appears as it does because the wave equation comes out to exactly that.\\

Also
\begin{align*}
	xu_x+yu_y=0 \implies u(x,y) = C
\end{align*}

\HLINE
\begin{description}
	\item \textbf{Observations}
	
	We have the various equations: Transport, Laplace/Poisson (basic), Heat and Wave equations.  These each have general solutions, usually depicted by $\Phi$ with specific solutions based on additional information.  The PDE is the entire equation and surrounding information: initial conditions, inhomogenous results, boundary conditions.  Further, some solutions are dependent on the 'architecture/structure' of the space which they are defined (e.g., Green's Functions).  Some special cases happen requiring the isolation of "spiked points" (points that go to infinity making the PDE non-harmonic).  These include surrounding these points with a boundary exclusion zone and integrating over the smooth subspace.
	
	There are additional clever tricks that must be considered.  Integration by parts using The Divergence Theorem (building a vector space where the dot product would be useful in this).  Dirac Functions.  Green's Functions and Green's Identities.
	
	Each equation has a Maximum Value Theorem, a Mean Value Theorem, Uniqueness Theorem.

	\item \textbf{Transportation Problem and Method Characteristics}
	
	Because the Transportation Equation is
	\begin{align*}
		u_t+b\cdot Du=0
	\end{align*}resembles a linear equation we can use that as inspiration.  Use
	\begin{align*}
		z(s) = u(x+sb, t+s)
	\end{align*}we'll see that $z' = 0$ which is a constant.  Thus, when $s=-t$, intial conditions, we get $z(s) = u(x-tb, 0)$ which is our intial contioins on $u$, namely $g(x)$.  Thus, $u(x,t) = g(x-tb)$
	
	\textbf{Non-Homogenous case}: We find that $z'(s) \ne 0$ and actually is equal to $z'(s) = f(x+sb, t+s$
	\begin{align*}
		u(x,t)-g(x-bt) &= z(0)-z(-t)\\
		&= \int_{-t}^1 z'(s)\, ds \\
		&= \int_0^t f(x+(s-tb,s)\,ds
	\end{align*}thus
	\begin{align*}
		u(x,t) &= g(x-bt) +\int_0^t f(x+(s-tb),s)\,ds
	\end{align*}
	
	\item \textbf{General Laplace solution.}
\begin{remark}[\textbf{Properties of $\Phi$}].

\begin{align*}
	\Phi(x) &= \Large\BINDEF{-\frac{1}{2\pi}\log |x| & (n=2)}{\frac{1}{n(n-2)\alpha(n)}\frac{1}{|x|^{n-2}} & (n\ge 3)} 
\end{align*}
\begin{align}
	D\Phi(y) &= \frac{-1}{n\alpha(n)}\frac{y}{|y|^n},\, (y\ne 0)\\
	\PART{\Phi}{\nu} &= \nu\cdot D\Phi(y)= \frac{-1}{n\alpha(n)\varepsilon^{n-1}},\, y \in \partial B(0,\varepsilon)  \\
	-\Delta\Phi &= \delta_0 \implies \Delta\Phi(x-y)=0,\, y \in U \AND y \ne x
\end{align}implying that there is a point source at 0 or $x$. In (2) $n\alpha(n)\varepsilon^{n-1}$ is the area of the surface of $\partial B(0,\epsilon)$

\end{remark}

	
	\item \textbf{General Poisson solution.}
	
	\begin{align*}
		&\BINDEF{-\Delta u =f, & x \in U}{u=g, & x \in \partial U}\\
		u(x,t) &= \Phi(x) - \phi^x(x) 
	\end{align*}where $\phi^x$ is a Dirichlet Delta function.
	
	\item \textbf{Green's Functions and Dirac Delta functions.}
	
	A Green's Function must satisfy the following:
	\begin{align*}
		G(x,y) &= \BINDEF{ \Delta G(x) = 0, & x \in U}{G(x,y) = \Phi(x-y) + \phi^x(x), & x,y \in \partial U}
	\end{align*}where $\phi^x$ is a Dirac Delta function determined by the geometry of $U$.  Then,
	\begin{align*}
		u(x) &= \int_U f(x)G(x,y)\, dy +\int_{\partial U} g(x)\PART{G}{\nu}(x,y) dS
	\end{align*}
	
	\item \textbf{Green's identities.}
\begin{remark}[Green's Identities].
\begin{align*}
\int_U \Delta u\; dx &= \int_{\partial U} \PART{u}{\nu} dS\\
	\int_U Dv\cdot Du\;  dx &= -\int_U u\Delta v \; dx + \int_{\partial U} \PART{v}{\nu}u\;  dS\\
	\int_U u\Delta v - v \Delta u\;  dx &= \int_{\partial U} u\PART{u}{\nu}-v\PART{u}{\nu}\;  dS
\end{align*}
\end{remark}
	
	\item \textbf{Neuman conditions.}
	
	Example from Practice Test 2:
	
	\item \textbf{Energy methods}
	
	Use the functional 
	\begin{align*}
		I[w] &= \int_U \frac{1}{2}|Dw|^2-wf\, dx \\
		w \in \mathcal{A} &= \{ w \in C^2\,|\, w=g \text{ on } \partial U\}
	\end{align*}Dirichlet's Principle then states that
	\begin{align*}
		I[u] &= \min_{w \in \mathcal{A}} I[w] \iff u \text{ is a solution}
	\end{align*}
	
	\item \textbf{Maximum Theorem}
	
	\begin{align*}
		\max_{\bar U} u &= \max_{\partial U} u
	\end{align*}
	
	\item \textbf{Mean Value Theorem}
	
	\begin{align*}
		u(x) &= \fint_U u\, dx = \fint_{\partial U} u \,dS
	\end{align*}
\end{description}

\HLINE
\begin{remark}[Integration by Parts\footnote{see page 9, 9\_17.pdf}].
\begin{align*}
	\int_U f\Delta g\, dx = -\int_U \nabla f \cdot \nabla g \, dx + \int_{\partial U} f \PART{g}{\nu}\, dS
\end{align*}
\begin{align*}
	-\Delta w &= 0 \implies w\Delta w = 0\\
	\int_U w\Delta w\; dx &= -\int_U |Dw|^2\;  dx+\int_{\partial U}wDw\cdot \nu\;  dx = 0\\
	\int_{\partial U} wDw\cdot \nu\;  dx &= \int_U |Dw|^2\;  dx + \int_U w\Delta w\;  \not dx
\end{align*}
\end{remark}

\newpage
\textbf{ChatGPT: Homework 3, question 5a: }

\begin{enumerate}[label=5(a)]
	\item Consider the initial value problem 
	\begin{align*}
		\BINDEF{u_t-\Delta u+b\cdot Du=0, & x\in \R^n, t>0}{u(x,0)=g(x) }
	\end{align*}where $b \in \R^n$ is a fixed vector. (Physically, the $b\cdot  Du$ term represents advection.)
	
	Find an explicit soluton to this probem using the  change of variables $z=x+bt$ and the solution formula for the heat equat1ion (equation (9) in Section 2.3 of Evans).
	
\end{enumerate}

\HLINE \\
Nice — this is the standard trick. The convenient choice is
$$
z := x - bt \qquad\bigl(\text{so }x=z+bt\bigr),
$$
and define
$$
v(z,t):=u(z+bt,t).
$$
Compute derivatives (keeping $z$ fixed when differentiating in $t$):
$$
v_t = u_t + b\cdot\nabla_x u,\qquad \nabla_z v=\nabla_x u,\qquad \Delta_z v=\Delta_x u.
$$
Substituting into the PDE $u_t-\Delta_x u+b\cdot\nabla_x u=0$ gives
$$
(v_t - b\cdot\nabla_x u) - \Delta_x u + b\cdot\nabla_x u = v_t-\Delta_z v =0.
$$
So $v$ satisfies the heat equation $v_t-\Delta_z v=0$ with initial data $v(z,0)=u(z,0)=g(z)$. Hence by the heat-kernel formula
$$
v(z,t)=\int_{\mathbb{R}^n}\Phi(z-y,t),g(y),dy,\qquad
\Phi(w,t)=\frac{1}{(4\pi t)^{n/2}}e^{-|w|^2/(4t)}.
$$
Returning to $u$ (recall $u(x,t)=v(x-bt,t)$) gives the explicit solution
$$
\boxed{u(x,t)=\int_{\mathbb{R}^n}\Phi(x-bt-y,t )g(y)\,dy
= \frac{1}{(4\pi t)^{n/2}} \int_{\mathbb{R}^n}
e^{-\frac{|x-y-bt|^2}{4t}}g(y)\,dy. }
$$

One can easily check this satisfies the PDE and $u(\cdot,t)\to g$ as $t\downarrow 0$. (If you prefer the change of variables $z=x+bt$ from your prompt, it is equivalent up to the sign of (b) — the above is the usual sign that removes the advection term.)

\newpage
In PDE terminology those two phrases refer to very different concepts:

---
\begin{description}

\item \textbf{ 1. **Bounded Domain**}

A *domain* is just the region in which the PDE is being considered.

* A **bounded domain** means this region is *finite in extent*.
* Example: The unit disk $ {x^2 + y^2 < 1} \subset \mathbb{R}^2 $ or the interval $ (0,1) \subset \mathbb{R} $.

This term is purely **geometric** — it says nothing about time or the type of PDE.

---

\item  \textbf{2. **Parabolic Boundary**}

This term is specific to **evolution PDEs**, especially **parabolic equations** like the heat equation.

Suppose you solve a PDE on a space-time cylinder:

$$
\Omega \times (0,T),\quad \text{where }\Omega\subset\mathbb{R}^n\text{ is the spatial domain}.
$$

The **parabolic boundary** $ \partial_p(\Omega \times (0,T)) $ consists of:

* The **initial time boundary**: $ \Omega \times {0} $
* The **lateral spatial boundary over time**: $ \partial\Omega \times (0,T) $

So for the heat equation, you don’t impose boundary conditions on the *final time* $t=T$, because information flows *forward* in time.

---

\item  \textbf{Summary Table}

\begin{tabular}{|l|l|l|l|}
\hline
Term  & Applies To  & Meaning  & Example \\
\hline
**Bounded domain**  & Any PDE  & Spatial region is finite & $ \Omega = (0,1) \subset \mathbb{R} $ \\
**Parabolic boundary** & Parabolic PDEs (e.g. heat equation) & Where boundary/initial data & $ \Omega\times{0} \cup \partial\Omega \times (0,T) $ \\
& &must be given in space-time & \\
\hline
\end{tabular}

\end{description}

Let me know if you'd like a sketch or diagram — this is much easier to see visually!\footnote{stuff and nonsense}{

\newpage
\textbf{\Large {October 29, 2025}}

\begin{description}
	\item Midterm \#2 Wednesday 11/5
	
	Heat Equation, Wave Equations, Fourier Transform.
	
	Practice test solutions and Homework \#4 solutions will be posted soon.
	
	\item More on Method of Characteristics (converting to a single variable ODE).
	
	\item \textbf{Conservation laws}. One dimension $ u_t + F(u)_x = 0 $, $u(x,t), x \in \R$.
	
	Higher Dimension: $u_t + \operatorname{div} [F(u)] = 0$ where $F$ is a flux and $u$ is a density.
	\begin{align*}
		\frac{d}{dt}\int_V u\, dx = \int_{\partial V} F(u)\cdot \nu\, dS
	\end{align*}The change is dependent on the flow over the boundary.
	
	Burger's Equation: $u_t + uu_x=0$ a non-linear PDE.
	
	Using Method of Characteristics
	\begin{align*}
		z(s) &= u(x(s), t(s)) \\
		\frac{dt}{ds} &= 1, \, \frac{dx}{ds} = u
	\end{align*}
	
	Rarefaction Wave.
	
	\item \textbf{Weak Formulation}
	
	Multiply PDE by a 'test function' $\phi$ that is smooth and compactly supported.
	\begin{align*}
		\int_0^\infty\int_{-\infty}^\infty \phi (u_t+uu_x)\, dx \, dt=0
	\end{align*}boundary terms aren't necessary because of compact support (closed and bounded domain).
	
	One benefit, the weak solution makes sense even when $u$ is not differenitable.
	
	This weak formulation helps us choose where find the interace curve.
\end{description}

\newpage
\textbf{\Large{November 3, 2025}}

\begin{description}
	\item \textbf{Test on Wednesday 11/5:} Review for today.
	
	\item 5 or 6 questions with formulas,   Not unlike the practice test and the homework assignments.
	
	\item \textbf{Heat Equation: }
	\begin{align*}
		u_t - u_{xx} &=0, \, x \in \R^n, t>0 \\
		\text{ or }	u_t - u_{xx} &=f, \, x \in \R^n, t>0 \\
		u(x,0) &= g(x), \text{ initial conditions} \\
		\text{Heat Kernel } \Phi(x,t) &= \frac{1}{(4\pi t)^{n/2}} e^{-|x|^2/4t} \\
		\int_{\R^n} \Phi(x,t) dx &= 1\end{align*}Derivation; use the scaling law ($\Phi(rx, r^2t)$ solves the same PDE), radial symmetry assumption, to get an ODE.  We solve this by taking the convolution (in $x$ only) 
		\begin{align*}
			u(x,t) = \Phi *_x g = \int_{\R^n} \Phi(x-y,t)g(y) dy
		\end{align*}We get a couple of things
		\begin{itemize}
			\item Instantaneous propogation $x$ is independent of $t$.
			\item Strict positive: $g \ge 0$ and $g \not \equiv 0$.
			\item Regularity: smoothness of solutions.
			\begin{align*}
				u(x,t) &= \int_{\R^n} \Phi(t, x-y) g(y)dy & (*)
			\end{align*}$\Phi(t, x-y)$ is smooth except at (0,0) (i.e., $t=0, x=y$).  Notice that 
		\begin{align*}
			\lim_{t\to 0} \frac{1}{(4\pi t)^{n/2}} e^{-|x|^2/4t} &= 0
		\end{align*}Hence the only problem that we have with (*) is when $x=y$ which can be resolved through the convolution.
		\end{itemize}
		
	\item \textbf{Non-homogeneous Heat Equation}
	
	\begin{itemize}
		\item Duhamel: 
		\begin{align*}
			u(x,t) &= \Phi *_{t,x} f = \int_0^t\int_{\R^n} \Phi(x-y, t-s) f(y,s)\, dy \, ds
		\end{align*} check text on adding a parameter $s$.
		
		\item Parabolic Maximum Principle 
		\begin{align*}
			\partial U \times [0, T) \cup ( U\times \{t=0\} )
		\end{align*}  doesn't include $t=T$ boundary!  Unique maximum should not occur at the end of time.
		
		\textbf{Theorem: } \begin{align*}
			\max_{\CONJ{U}_T} \le \max_{\Gamma U_t}
		\end{align*}
		\item Uniqueness derived from the Maximum Principle.  Therefore (*) is THE solution.
		
		\item Regularity on $\R^n$ for $u_t - \Delta_{xx}=0$: 
		\begin{itemize}
			\item On $\R^n$, use the solution formula.
			\item On a bounded domain: use a localizing cutoff.  Cylinder proof
		\end{itemize}
		
		\item \textbf{Energy Methods} (too much done that it seems to be the entire course).  Two ways in: 
		\begin{itemize}
		\item given some functional, $\int u(s)\,dx$ differentiate, use PDE, integrate by parts (almost always used), then try to show $t$-derivative of the functional has a sign.
		
		\item or mulitply PDE by something, integrate by parts, find the t-derivative of some functional that is constant or decreasing.
		
		\end{itemize}
	\end{itemize}
	
	\item \textbf{the Wave Equation:} the preponderance of questions will be in one-dimension.  Just to make it worthy of a decent test question.  Starting d'Alembert formula.
	\begin{align*}
		(\partial_t+\partial_x)(\partial_t-\partial_x)u = 0
	\end{align*}a transport equation. 
	
	\begin{itemize}
		\item Finite speed of propagation.  See second term in d'Alembert
		\item Use the method of spherical means
		\item 2D method of descent (go to 3 and ignore one coordinate).
		\begin{align*}
			u_{tt}-u_{xx}- u_{yy} = 0
         \end{align*}	Huygens Principle.
         
         \item non-homogenous: use Duhamel
         
		
		\BLUE{\item \textbf{Energy methods} are good for an integral estimate.  Maximum methods are good for point-wise.\\
		First attempt\begin{align*}
			\int_0^l u_{tt}\,u_t \,dx &= \int_0^l c^2u_{xx}\,u_t \,dx \\
			\frac{1}{2} \frac{d}{dt}\int_0^l (u_t)^2\, dx &= c^2\PAREN{u_x(l,t)u_t(l,t)-u_x(0,t)u_t(0,t)-\int_0^l u_x u_{tx} \,dx} \\
			&= -c^2 \frac{1}{2}\frac{d}{dx} \int_0^l (u_x)^2 dx
		\end{align*}where $u(0,t) = 0 \implies u_t(0,t)=u_t(l,t) = 0$.  Therefore, define your energy function $E$
		\begin{align*}
			E(t) &= \frac{1}{2}\PAREN{\int_0^l (u_t)^2 + c^2 (u_x)^2\, dx} \\
			\text{notice } E'(t) &= 0 \implies E(t) \equiv c
		\end{align*}Hence a conservation of energy.  Further $(u_t)^2$ is the kinectic energy and $c(u_x)^2$ is the potential energy and the wave passes this energy back and forth smoothly.  This leads to uniqueness of the wave equation.
		}
		\item \textbf{No maximum principle}
	\end{itemize}
	
	\item \textbf{Fourier Transform:}  Take a function 
	\begin{align*}
	\mathcal{F}[f] &= \hat{f}(w) = \int_\R e^{-i\omega x}f(x) dx \AND \mathcal{F}^{-1}[f] = \hat{f}(\omega ) = \int_\R e^{-i\omega x}\hat{f}(x) dx \\
		\mathcal{F}['f(x)] &=  \int_\R e^{-i\omega x}f'(x) dx = i\omega  \hat{f}(\omega )
	\end{align*}

	\textbf{General Idea: Heat Equation}
	\begin{enumerate}
		\item \textbf{Step One: }Because integration is from $-\infty$ to $\infty$.  We begin by converting to Fourier Transform
		\begin{align*}
			u_t &= \alpha u_{xx} \\
			\FFF(\hat{u}_t) &= \FFF(\alpha u_{xx})= -\alpha^2 \omega \hat{u} \\
		\end{align*}
		\item Solve this ODE with respect to $\omega$
		\begin{align*}
			\int \frac{\partial \hat{u}}{\hat{u}} &= \int -\alpha^2 \omega^2\, \partial t \\
			\hat{u}(\omega, t) &= \varphi(\omega)e^{-\alpha^2\omega^2 t}
		\end{align*}for some initial condition $\varphi(\omega)$
		\item \textbf{Inverse the Fourier Transform} and take use the convolution identity $\FFF^{-1}(f * g) = \FFF^{-1}(f)\FFF(g)$
		\begin{align*}
			u(x,t) &= \FFF^{-1}(\varphi(\omega)e^{-\alpha^2\omega^2 t}) \\
			&= \FFF^{-1}(\FFF(\varphi(\omega))\FFF(e^{-\alpha^2\omega^2 t})) \\
			&= \FFF^{-1}\PAREN{(\varphi(\omega))*(e^{-\alpha^2\omega^2 t})} \\
			&= \varphi(\omega)*(e^{-\alpha^2\omega^2 t})
		\end{align*}
		
		\item \textbf{In summary: }  These identities make the use of Fourier Transforms possible
		\begin{align*}
			\frac{d^n u}{dx^n} &= ( i\omega)^n u \\
			\FFF(f * g) &= \FFF(f)\FFF(g) \\
			\FFF(\FFF^{-1}(u)) &= u
		\end{align*}
	\end{enumerate}
	
	\BLUE{\item \textbf{Strategies:}
	\begin{itemize}
		\item \textbf{Solutions based on Integration by Parts for integrals}
		\item \textbf{Solutions based on Maximum Principle for point-wise.}
	\end{itemize}
	}
\end{description}

\newpage
\textbf{\Large{November 12, 2025}}

\textbf{Test returned:} Average was on the "low side".  Curve to date of cumumlative grade (unofficial):\\
\begin{tabular}{ll}
	 A & 65\% and up,\\
	 B & 50\% - 64\%,\\
	 C & 49\% or less.
\end{tabular}
\\

According Canvas I'm at 49.81\%!  Sacka fracka gall durn it!  I've got to study to get a B!  Final exam is 30\% of grade $50\% * 70\% = 35\%.$  if I get 50\% on the final I might squeak by with a B-.  I could get an A by getting a perfect score, which I sincerely doubt.  But, that being the case, there is a chance.  Use the Dr. Peyham videos on YouTube to study.  (28 lessons and 28 days!)\\

\textbf{Questions from the test.}
\begin{enumerate}
	\item Show solutions to $u_{tt}-u_{xx}=0$ do NOT necessarily satisfy the max condition.  Real number line only.  Proof by counter example.  Let $g \equiv 0$ and $h \equiv 1$.
	
	\item Differentiate in time and see what happens then plug and play, follow by Integrate by Parts (which no one seemed to get right).
	
	\item (skipped)
	
	\item PDE like the wave equation but not.  Imitating the derivation of d'Alembert's Formula.
	
	\item (skipped)
\end{enumerate}

My observation:  Today he went over the test and, frankly, it looks very familiar.  He gave me 30\%, though, so I have to wonder what I actually did wrong.  There was one problem that made no sense to me because it required some knowledge about Statistics.  I forgot how to solve the Transport Equation which probably accounted for a lot of lost points.  \textbf{\textit{Still dozing off.}}

\HLINE
\textbf{\Large{Weakly harmonic solutions are harmonic to PDE}}

\newpage
\textbf{\Large{November 17, 2025}}\\

Wednesday's class will be ``asynchronous" that can be found on Panopto in Canvas.  Homework \#5 will be posted today and due on December 1, 2025.\\

Continuing on Weakly Harmonic Functions and their connection to Calculus of Variations.\\

$u \in C(U)$ is \textit{weakly harmonic} if $\forall$ smooth test functions $\phi$ with compact support in $U$ 
\begin{align*}
	\int_U u\Delta \phi\,dx = 0
\end{align*}If $\Delta u=0 \implies u$ is weakly harmonic $0=\int_u\phi\Delta u \, du = \cdots \int u\delta\phi\, dx$\\

Convergence of Dirac Energy Expression: does it converge? We need Functional Analysis, using compactness theorems in function space.

\HLINE
Special sort of PDE using a matrix  div$ (A(u)Du) = 0$.  Which becomes the Laplacian when $A = I$.  The PDE is \textit{elliptic} if $A(x)$ is positive definite.  If $A$ is symetric and positive definite then its eigenvalues, all its eigenvalues are positive.  ``Elliptic" means like the Laplace Equation (how most people think of it).  And \textit{uniformaly elliptic} when all eigenvalues are contained in an interval $[\lambda, \Lambda]$.\\

Anisotropic difusion.\\

\HLINE
Classification of 2$\SND$-order linear PDEs.  The general form in 2 variables.
\begin{align*}
	Au_{xx}+Bu_{xy}+Cu_{yx}+Du_{x}+Eu_y+Fu=0
\end{align*}where not all $A,B,C,D$ are zero.  Hence we have
\begin{align*}
	B^2-4AC &<0 & \text{ellpitical, Laplace} \\
	B^2-4AC &>0 & \text{hyperbolic, Heat} \\
	B^2-4AC &=0 & \text{parabolic, Wave} 
\end{align*}which are named through analogy.  Hence Laplace is elliptic, the Heat Equation is parabolic, the Wave Eqation is hyperbolic.  In practice, $A,B,C$ do NOT take the form of functions.\\

Classifications of higher dimensional 2\SND order PDEs.

\newpage
\noindent\textbf{\Large{New Classifications for solving PDEs}}\\

Given the following classificatin of 2$\SND$-order Partial Differential Equations
\begin{align*}
	Au_{xx}+Bu_{xy}+Cu_{yx}+Du_{x}+Eu_y+Fu=0
	\end{align*}If $A=B=C=0$ and $F=0$ we have the \textbf{transport equation}.  Then, where not all $A,B,C,D$ are zero.  Hence we have
\begin{align*}
	B^2-4AC &<0 & \text{ellpitical, Laplace} \\
	B^2-4AC &>0 & \text{hyperbolic, Heat} \\
	B^2-4AC &=0 & \text{parabolic, Wave} 
\end{align*}

\begin{itemize}

\item Classifications
\begin{itemize}
	\item Transport Equation
	\item Elliptical
	\item Hyperbolic
	\item Parabolic
	\item Burgers' Equation (Misc?)
\end{itemize}

\item Then we define sub-classifications/features of PDE as
\begin{itemize}
	\item Homogenous vs inhomogenous.
	\item Boundary conditions.
	\item Initial conditions.
\end{itemize}

\item And Strategies for solving PDEs
\begin{itemize}
	\item Method of Characteristics
	
	By finding a contour upon which all values of $u(x,t)$ are the same, we can then use the slope of that contour (by definition a constant) as part of the PDE, reducing it to an ODE.
	
	\item Integration by parts (in order to use the Divergence Theorem).
	
	\begin{align*}
		\int_U w\Delta w\; dx &= -\int_U |Dw|^2\;  dx+\int_{\partial U}wDw\cdot \nu\;  dx = 0 \\
	\int_U Dv\cdot Du\;  dx &= -\int_U u\Delta v \; dx + \int_{\partial U} \PART{v}{\nu}u\;  dS
	\end{align*}
	\item Energy methods.
\end{itemize}

\end{itemize}

\noindent\textbf{Theoretical Constructs}
\begin{itemize}
	\item Max/Min Principle
	\item Divergence Theorem
	\item Average integral usage.
\end{itemize}

\noindent\textbf{Solutions via Classification: }We find solutions to PDEs via their classifications, all of which have are meant to reduce a PDE to an ODE.

\begin{description}
	\item \textbf{Advection/Transport Equation: $u_t+au_x=0$}
	\begin{enumerate}
		\item Method of Chacteristics
		
		\BOXIT{$u(s)$ must be constant satisfying
		\begin{align*}
			\frac{dx}{dt} &= a \\
			x(t) &= x_0+at\\
			u(x,t) &= u_0(x-at)
		\end{align*}}
		
		\item Fouirer Transform
		
		\begin{align*}
			\hat{u}_t + ia\xi\hat{u}=0
		\end{align*}solves the ODE to
		\begin{align*}
			\hat{u}(\xi,t)&= u_0(\xi)e^{ia\xi t}
		\end{align*}then using the inverse Fourier Transform
		\begin{align*}
			u(x,t) &= u_0(x-at)
		\end{align*}
	\end{enumerate}
	\item \textbf{Elliptical: Laplace/Poison, $u_t-\Delta u=0$}
	\item \textbf{Hyperbolic: The Heat Equation}
	\item \textbf{Parabolic: The Wave Equation, $u_{tt}-c^2u_{xx}=0$}
	\begin{itemize}
		\item With Initial Conditions'
		\begin{align*}
			\BINDEF{u_{tt}-c^2u_{xx}=0, & t>0, x\in \R^n}{u(x,0) = f(x), \\ u_t(x,0) = g(x)}
		\end{align*}use d'Alembert
		\begin{align*}
			u(x,t) &= \frac{1}{2}\PAREN{f(x+ct)+f(x-ct)}+\frac{1}{2}\int_{x-ct}^{x+ct} g(s)ds
		\end{align*}
	\end{itemize}
\end{description}
\newpage
\textbf{\Large{Greens's Function}}

The request asks for a valid Green's function used in Partial Differential Equations (PDEs). A Green's function is a fundamental solution or impulse response for a linear, non-homogeneous differential operator defined on a domain, subject to specified boundary conditions.\\

\textbf{ Reasoning \& Explanation}

\begin{description}

\item \textbf{1. Definition of a Green's Function}

A Green's function, $G(\mathbf{x}, \mathbf{x}')$, for a linear differential operator $L$ acting on a function $u(\mathbf{x})$ in a domain $\Omega$, is defined as the solution to the equation:
$$L G(\mathbf{x}, \mathbf{x}') = \delta(\mathbf{x} - \mathbf{x}')$$
where $\delta(\mathbf{x} - \mathbf{x}')$ is the Dirac delta function, representing a point source located at $\mathbf{x}'$.

If the Green's function is known, the solution $u(\mathbf{x})$ to the non-homogeneous equation $L u(\mathbf{x}) = f(\mathbf{x})$ (with appropriate boundary conditions) can be found by integrating the source function $f(\mathbf{x})$ against the Green's function:
$$u(\mathbf{x}) = \int_{\Omega} G(\mathbf{x}, \mathbf{x}') f(\mathbf{x}') d\mathbf{x}'$$

\item \textbf{2. Example: Green's Function for the Laplacian Operator (Poisson Equation)}

The most common and fundamental example of a Green's function in PDEs is associated with the Laplacian operator, $\nabla^2$. This operator appears in the Poisson equation ($\nabla^2 u = f$) and Laplace's equation ($\nabla^2 u = 0$).

The defining equation for the Green's function $G(\mathbf{x}, \mathbf{x}')$ in this context is:
$$\nabla^2 G(\mathbf{x}, \mathbf{x}') = \delta(\mathbf{x} - \mathbf{x}')$$

We consider the Green's function in different dimensions:

\begin{description}

\item  \textbf{A. 3-Dimensional Space ($\mathbb{R}^3$)}

In three dimensions, $\mathbf{x} = (x_1, x_2, x_3)$, and the distance $r = |\mathbf{x} - \mathbf{x}'|$. The fundamental solution (Green's function for the unbounded domain) is:
$$G_3(\mathbf{x}, \mathbf{x}') = -\frac{1}{4\pi r} = -\frac{1}{4\pi |\mathbf{x} - \mathbf{x}'|}$$

**Verification (Conceptual):** The Laplacian of $1/r$ is proportional to the Dirac delta function. Specifically, $\nabla^2 \left(\frac{1}{r}\right) = -4\pi \delta(\mathbf{x} - \mathbf{x}')$. Therefore, multiplying by $-1/(4\pi)$ yields the required Green's function.

\item \textbf{B. 2-Dimensional Space ($\mathbb{R}^2$)}

In two dimensions, $\mathbf{x} = (x_1, x_2)$, and $r = |\mathbf{x} - \mathbf{x}'|$. The fundamental solution is:
$$G_2(\mathbf{x}, \mathbf{x}') = \frac{1}{2\pi} \ln(r) = \frac{1}{2\pi} \ln(|\mathbf{x} - \mathbf{x}'|)$$

**Verification (Conceptual):** The Laplacian of $\ln(r)$ is proportional to the Dirac delta function. Specifically, $\nabla^2 (\ln r) = 2\pi \delta(\mathbf{x} - \mathbf{x}')$. Therefore, multiplying by $1/(2\pi)$ yields the required Green's function.

\end{description}

\item \textbf{3. Conclusion}

A valid Green's function used in PDEs is the fundamental solution for the Laplacian operator in 3D space.

$$G(\mathbf{x}, \mathbf{x}') = -\frac{1}{4\pi |\mathbf{x} - \mathbf{x}'|}$$

This function is crucial in solving the Poisson equation in electrostatics (potential due to a charge distribution) and gravitation (potential due to a mass distribution).

\item \textbf{Final Answer Summary}

A valid Green's function used in Partial Differential Equations (PDEs) is the fundamental solution for the **Laplacian operator ($\nabla^2$)** in three-dimensional Euclidean space ($\mathbb{R}^3$).

%| Operator | Defining Equation | Domain | Valid Green's Function $G(\mathbf{x}, \mathbf{x}')$ |
%|---|---|---|---|
%| Laplacian ($\nabla^2$) | $\nabla^2 G = \delta(\mathbf{x} - \mathbf{x}')$ | $\mathbb{R}^3$ (Unbounded) | $G(\mathbf{x}, \mathbf{x}') = -\frac{1}{4\pi |\mathbf{x} - \mathbf{x}'|}$ |
%| Laplacian ($\nabla^2$) | $\nabla^2 G = \delta(\mathbf{x} - \mathbf{x}')$ | $\mathbb{R}^2$ (Unbounded) | $G(\mathbf{x}, \mathbf{x}') = \frac{1}{2\pi} \ln(|\mathbf{x} - \mathbf{x}'|)$ |

Final Answer:
\begin{enumerate}
\item \textbf{A valid Green's function used in PDEs is the fundamental solution for the Laplacian operator in three-dimensional space: \(G(\mathbf{x}, \mathbf{x}') = -\frac{1}{4\pi |\mathbf{x} - \mathbf{x}'|}\)}

\item Identifying a valid Green's function used in PDEs
Definition of a Green's function
\item A Green's function \(G(\mathbf{x}, \mathbf{x}')\) for a linear differential operator \(L\) satisfies \(L G(\mathbf{x}, \mathbf{x}') = \delta(\mathbf{x} - \mathbf{x}')\), where \(\delta\) is the Dirac delta function.
Fundamental solution for the Laplacian in 3D
\item The Green's function for the Laplacian \(\nabla^2\) in unbounded three-dimensional space is \(G(\mathbf{x}, \mathbf{x}') = -\frac{1}{4\pi |\mathbf{x} - \mathbf{x}'|}\). This satisfies \(\nabla^2 G(\mathbf{x}, \mathbf{x}') = \delta(\mathbf{x} - \mathbf{x}')\) and is widely used in solving Poisson's equation.

\end{enumerate}
\textbf{Answer:} \textbf{A valid Green's function used in PDEs is the fundamental solution for the Laplacian operator in three-dimensional space: \(G(\mathbf{x}, \mathbf{x}') = -\frac{1}{4\pi |\mathbf{x} - \mathbf{x}'|}\)}

\end{description}

\newpage
\section*{More on Greens Functions}

\textbf{Green’s Functions for PDEs — Intuition First}

\DEFINE{Green’s functions} are a method for solving linear differential equations (especially PDEs) by breaking the problem into responses to **point sources** and then adding them up.

Think of it like this:

> If you know how a system reacts to a single impulse, you can build the solution for any input by adding impulses everywhere.

This idea appears everywhere: heat flow, electrostatics, waves, quantum mechanics.

---

\begin{description}

\item 1. The Core Idea

Suppose we have a linear PDE:
\begin{align*}
L u(x) = f(x)
\end{align*} Where:
\begin{itemize}
\item $L$ = differential operator (like Laplacian)
\item $u(x$) = solution
\item $f(x$) = forcing/source term
\end{itemize}
Instead of solving this directly, we solve a simpler problem:
\begin{align*}
L G(x,\xi) = \delta(x-\xi)
\end{align*}
This function \DEFINE{G} is the \DEFINE{Green’s function}.

It represents the response at point $x$ due to a point source located at $\xi$.

---

\item 2. Once You Have the Green’s Function

The solution to the original PDE becomes:
\begin{align*}
u(x) = \int G(x,\xi) f(\xi), d\xi
\end{align*}

\DEFINE{Interpretation:}
\begin{itemize}
\item Every point $ \xi $ acts like a tiny source
\item $G(x,\xi$) tells how much that source affects point $x$
\item Add all contributions → solution
\end{itemize}

This is basically \DEFINE{continuous superposition}.

---

\item 3. Physical Intuition Examples

\item Heat Equation

Green’s function = temperature spread from a \DEFINE{single heat pulse}
\begin{itemize}
\item Electrostatics (Poisson equation)

Green’s function = potential from a \DEFINE{point charge}

\item Wave equation
\end{itemize}

Green’s function = wave produced by a \DEFINE{single impulse}

So:

\begin{tabular}{|l|l|}
\hline
 PDE       & Green’s function meaning    \\
\hline
 Heat      & Spread of heat from point   \\
 Poisson   & Potential from point charge \\
 Wave      & Wave from impulse           \\
 Diffusion & Spread from particle        \\
\hline
\end{tabular}

---

\item 4. Example: 1D Poisson Equation

Solve:
\begin{align*}
-\frac{d^2 u}{dx^2} = f(x)
\end{align*}

Green’s function satisfies:
\begin{align*}
-\frac{d^2 G}{dx^2} = \delta(x-\xi)
\end{align*}

The solution ends up being piecewise linear:

\begin{align*}
G(x,\xi) = \BINDEF{ x(1-\xi), & x < \xi }{ \xi(1-x), & x > \xi}
\end{align*}

Then:
\begin{align*}
u(x) = \int_0^1 G(x,\xi) f(\xi), d\xi
\end{align*}

---

\item 5. Big Picture Concept (Very Important)

Green’s functions convert:

\item Differential equation → Integral equation

This is powerful because:

\begin{itemize}
\item Integrals are often easier
\item Works with boundary conditions
\item Used in physics everywhere
\item Foundation of FEM, BEM, quantum field theory
\end{itemize}

---

\item 6. The Deep Insight

Green’s function is essentially the \DEFINE{inverse of a differential operator}.

\begin{align*}
G = L^{-1}
\end{align*}

So solving:
\begin{align*}
Lu = f
\end{align*}

becomes:
\begin{align*}
u = L^{-1} f
\end{align*}

And the inverse operator is represented by the integral with $G$.

---

\item 7. If You Remember Only One Thing

\DEFINE{Green’s function = response of the system to a point source.}

\DEFINE{Solution = sum of responses to all point sources.}

---

\item 8. Mental Picture

Imagine dropping pebbles into a pond:

\begin{itemize}
\item One pebble → ripple = Green’s function
\item Many pebbles → total waves = solution
\item Add all ripples together
\end{itemize}

That’s literally Green’s functions.

---

\item 9. Where They Are Used

Major areas:

\begin{itemize}
\item Electrostatics
\item Heat transfer
\item Fluid dynamics
\item Quantum mechanics
\item Signal processing
\item Control theory
\item Finite element methods
\item Machine learning kernels

\end{itemize}
---

\item If you want, next I can show:

\begin{itemize}
\item How to derive a Green’s function
\item Green’s functions for Laplace equation
\item Heat equation Green’s function
\item Fourier transform connection
\item Eigenfunction expansion method
\end{itemize}
\end{description}


\newpage
\textbf{\Large{Mollifiers and Mollification}}

Mollifiers and the process of mollification are indeed related to the concepts discussed earlier, specifically concerning the Dirac delta function, which is central to the definition of a Green's function. While a Green's function is a solution to an equation with a delta function source, mollifiers are used to approximate functions, including distributions like the delta function, with smooth functions.

\begin{description}
\item \textbf{Mollifiers}

A mollifier is a special type of smooth function used to create smooth approximations of other functions, particularly those that are not smooth or are generalized functions (distributions). A standard mollifier, often denoted as \(\eta(x)\), has the following properties:

\begin{enumerate}

\item  **Smoothness:** 

\(\eta(x)\) is infinitely differentiable (i.e., \(\eta \in C^\infty\)).
\item **Compact Support:** 

\(\eta(x) = 0\) for \(|x| \ge 1\). This means it is non-zero only within a bounded region (the unit ball in \(\mathbb{R}^n\)).
\item **Normalization:** 

The integral of the function over its entire domain is 1: \(\int_{\mathbb{R}^n} \eta(x) dx = 1\).

\end{enumerate}
A common example of a mollifier in \(\mathbb{R}^n\) is:

\[ \eta(x) = \begin{cases} C \exp\left(\frac{1}{|x|^2 - 1}\right) & \text{if } |x| < 1 \\ 0 & \text{if } |x| \ge 1 \end{cases} \]

where the constant \(C\) is chosen to ensure the integral is 1.

\item \textbf{Mollification}

Mollification is the process of using a mollifier to smooth out a function. This is done via convolution. For a function \(f\), its mollified version \(f_\epsilon\) is defined as:

\[ f_\epsilon(x) = (f * \eta_\epsilon)(x) = \int_{\mathbb{R}^n} f(y) \eta_\epsilon(x-y) dy \]

Here, \(\eta_\epsilon(x) = \frac{1}{\epsilon^n} \eta\left(\frac{x}{\epsilon}\right)\) is a scaled version of the standard mollifier \(\eta(x)\). As the parameter \(\epsilon \to 0\), the support of \(\eta_\epsilon\) shrinks, and it becomes more concentrated around the origin, behaving like an approximation of the Dirac delta function. Consequently, the mollified function \(f_\epsilon\) converges to the original function \(f\) in an appropriate sense.

\item \textbf{ Application in Solving PDEs}

Mollification is a powerful theoretical tool in the study of PDEs, particularly for proving the existence and regularity of solutions:

\begin{enumerate}

\item **Approximation:** 

It allows us to approximate a non-smooth (weak) solution of a PDE with a sequence of smooth functions. We can then apply classical calculus and analysis techniques to these smooth approximations.
\item  **Existence Proofs:** 

In proving the existence of weak solutions, one might first solve a regularized or mollified version of the PDE. Then, by taking the limit as \(\epsilon \to 0\), one can show that the sequence of smooth solutions converges to a weak solution of the original PDE.
\item  **Regularity Theory:** 

Mollification is used to prove that weak solutions are actually smoother than initially assumed (regularity). For instance, one can show that if a function is a weak solution to an elliptic PDE, its mollified version also satisfies a related equation, which can then be used to derive estimates on the derivatives of the solution.

\end{enumerate}

\item \textbf{In summary,} just as the Dirac delta function is a key concept for defining Green's functions, mollifiers serve as smooth approximations to the delta function. This connection makes mollification a fundamental technique for rigorously establishing the properties of solutions to PDEs, bridging the gap between generalized functions and classical, smooth functions.
\end{description}


\newpage
\textbf{\Large{December 1, 2025} \\Final exam 12/8 1:00 - 3:00 in Crawford 332}

\begin{description}
	\item \textbf{Final:} \DEFINE{Cumulative.}
	
	Practice Final has been published.	
	
	Eigenvalue problems
	
	conservation laws
	
	shocks
	
	classification
	
	\item From the begining of the course

	Important tactics.
	\begin{itemize}
		\item Method of Characteristics (contours).
		
		Its really much simpler than they say.  
		\begin{align*}
			au_t+bu_x=c \implies \frac{dt}{ds}=a, \, \frac{dx}{ds}=b, \, \frac{du}{ds}=c
		\end{align*}Where $a,b,c$ can be constants or functions.  Then, anti-differentiate based on $s$ and solve removing $s$ to find a solution.
		
		\item Understanding and using the Kernel. Represented by $\Phi$ and depends on the classification.
		\item Energy methods.
		\item Convolution.
		\item Mollification.
		\item Green's Function.
		
		See previous section.
		
		\item Dirac Delta functions.
		\item \textbf{Fourier Transform Methods}  One dimensional case.  The analog for a Fouier Series.  Basically, transform the equation, solve the ODE and transform back.  Identities
	
	\begin{align*}
		\mathcal{F}\PAREN{\frac{du}{dx_j}} &= i\xi_j\hat{u}(\xi)
	\end{align*}where $\xi_j$ is the frequency component corresponding to the coordinate $x_j$.  Thus, the Fourier transform of a derivative is reduced to a multiplication.  Then
	\begin{align*}
		\mathcal{F}(\Delta u) &= -|\xi|^2\hat{u}(\xi)  =  \hat{f}(\xi)\\
		&\AND \\
		\hat{u}(\xi) &= -\frac{1}{|\xi|^2}\hat{f}(\xi)
	\end{align*} then transforming back (using $\mathcal{F}^{-1}(ab)=\mathcal{F}^{-1}(a)*\mathcal{F}^{-1}(b)$
	\begin{align*}
		u(x) &= \mathcal{F}^{-1}\PAREN{-\frac{1}{|\xi|^2}}* f(x)
	\end{align*}where the term $\mathcal{F}^{-1}\PAREN{-\frac{1}{|\xi|^2}}$ is the Green's Function from above.  Thus,  $$\mathcal{F}^{-1}\PAREN{-\frac{1}{|\xi|^2}}=\frac{1}{4\pi},\, n=3$$
	
	Identity of derivative.
	
	Convolution Rule. $\widehat{f*g} = \hat{f}\hat{g}$
	
	Don't expect much of this but check out the type of question on test 2.

		\item construct a space where the PDE makes sense with a dot-product.
	\end{itemize}		
	
	Ingteresting additional topics
	\begin{itemize}
		\item Mean Value Principles.
		\item Maximum value principles.
		\item Regularity
	\end{itemize}
	
	\textbf{Classifying the space of PDEs}:

Given the following classificatin of 2$\SND$-order Partial Differential Equations
\begin{align*}
	Au_{xx}+Bu_{xy}+Cu_{yx}+Du_{x}+Eu_y+Fu=0
	\end{align*}If $A=B=C=0$ and $F=0$ we have the \textbf{transport equation}.  Then, where not all $A,B,C,D$ are zero.  Hence we have
\begin{align*}
	B^2-4AC &<0 & \text{ellpitical, Laplace} \\
	B^2-4AC &>0 & \text{hyperbolic, Heat} \\
	B^2-4AC &=0 & \text{parabolic, Wave} 
\end{align*}

\begin{itemize}

\item Classifications
\begin{itemize}
	\item Transport Equation
	\item Elliptical
	\item Hyperbolic
	\item Parabolic
	\item Burgers' Equation (Misc?)
\end{itemize}

\item Then we define sub-classifications/features of PDE as
\begin{itemize}
	\item Homogenous vs inhomogenous.
	\item bounded and unbounded.
	\item Boundary conditions.
	\item Initial conditions.
\end{itemize}

\end{itemize}
	\textbf{Strategies for each classification}
	\begin{itemize}
		\item \textbf{Generic Transport $au_t+bu_x+ cu=f(x,t)$}
		
		Use Method of Characteristics (beware the bounded domain, would need boundary/initial condition.  Which do you hit first the $t=0$ of $x=0$?).
		
		\textit{Must know this in order to be able to do the more complex solutions below.}
		
		\textbf{Non-Homogenous case}: We find that $z'(s) \ne 0$ and actually is equal to $z'(s) = f(x+sb, t+s)$
	\begin{align*}
		u(x,t)-g(x-bt) &= z(0)-z(-t)\\
		&= \int_{-t}^1 z'(s)\, ds \\
		&= \int_0^t f(x+(s-tb,s)\,ds
	\end{align*}thus
	\begin{align*}
		u(x,t) &= g(x-bt) +\int_0^t f(x+(s-tb),s)\,ds
	\end{align*}
		
		\item \textbf{Elliptical: Laplace's/Poisson's Equation $\Delta u=0$, $-\Delta u=f$.}
		
		Beware the domain, bounded domain vs whole space.
		
		Fundamental solution in $\R^n$, $\Phi(x)$
		
\begin{remark}[\textbf{Properties of $\Phi$}].

\begin{align*}
	\Phi(x) &= \Large\BINDEF{-\frac{1}{2\pi}\log |x| & (n=2)}{\frac{1}{n(n-2)\alpha(n)}\frac{1}{|x|^{n-2}} & (n\ge 3)} 
\end{align*}
\begin{align}
	D\Phi(y) &= \frac{-1}{n\alpha(n)}\frac{y}{|y|^n},\, (y\ne 0)\\
	\PART{\Phi}{\nu} &= \nu\cdot D\Phi(y)= \frac{-1}{n\alpha(n)\varepsilon^{n-1}},\, y \in \partial B(0,\varepsilon)  \\
	-\Delta\Phi &= \delta_0 \implies \Delta\Phi(x-y)=0,\, y \in U \AND y \ne x
\end{align}implying that there is a point source at 0 or $x$. In (2) $n\alpha(n)\varepsilon^{n-1}$ is the area of the surface of $\partial B(0,\epsilon)$

\end{remark}
		Convolusion $\Phi * f$ solves Poisson's Equation.  Works if $f$ decays, helps make it well-defined, i.e., $f$ is bounded and compactly supported.  Doesn't work at the singularity. Use a Dirac Delta function (not likely to be on the final).
		\begin{itemize}
			\item \textbf{Mean Value Property}, i.e., harmonic implies mean on a ball is value of center.
			
			able to prove nice things about harmonic functions, i.e., $\Delta u=0$
			
			Can prove maximum principle
			
			\textbf{Regularity}: if $u \in C(U)$ satisfies MVP implies $u \in C^\infty(U)$.  Prove using standard mollificontion.  i.e., convolve $u$ then show that $u$ equals the standard mollification. 
			
			\begin{description}
				\item Outline of proof: MVP proves that the integral of a harmonic function over a ball is the value at its center.   Differentiate this, proving that the derivative is also harmonic, i.e., a member of $C^1$. ``Bootstrap" this process by differentiating again (rinse and repeat) proving that it is actually $C^\infty$.
			\end{description}
			
			\textbf{Quantitative Regularity}: estimate on the derivatives. Liouville's Theorem for Harmonic Functions ($u$ harmonic and bounded implies $u$ is a constant).
			
			\textbf{Harnack Inequality}: a quantitative approach to the ``topology" of the set $U$ overwhich $u$ is defined.  Put simply:
			\begin{align*}
				\exists C \to 0<u(x_1) \le Cu(x_2), \,\forall x_1, x_2 \in V \subseteq U
			\end{align*}
			
			\textbf{Greens Functions:} solving the problem on a bounded domain (as opposed to the whole space). Connect to $\Phi$ and some $\phi^x$ ``can get a Green's Formula somewhat reliably!" Don't have this Newton potential.  Generally, helps to have symmetry.
			
			Leads to Integral formulas for ball, half space.
			
			\textbf{Energy Methods:} multiply the PDE which allows for integration by parts.
			\begin{itemize}
				\item Proves Uniqueness can be proved.
				\item Proves linearity of PDE (subtract two solutions get another solution)
				
				\item Converse of MVP: that is, MVP implies harmonic.

			\end{itemize}
						
			\item \textbf{Dirichlet's Principle.}
			
			The solution for 
			\begin{align*}
				\BINDEF{\Delta u, & x \in U}{u(x)=g(x), & x \in \partial U}
			\end{align*}is the function that minimizes \textit{Dirichlet Energy}
			\begin{align*}
				E[v] = \int_U |\Delta u|^2\,dx
			\end{align*}over the set of admissable functtions on the boundary.
			
			\item \textbf{On the test} use some of these details.  stuff where you'd have to just ``know" to use MVP versus something else.
			
		\end{itemize}
		
		\item \textbf{Hyperbolic: Heat Equation $u_t-\Delta u=0$}
		
		\textbf{Scaling Symmetry: }If $u$ solves the Hyperbolic PDE so does $u_\lambda(x,t) = u(\lambda x, \lambda^2 t)$, because of the chain rule. Use convolution and \textbf{The Heat Kernel, $\Phi$} (not the same as above)
		
		Nicer regularity properties than the Newton Potential.
		
		Infinite speed of propagation (not true in other PDEs).
		
		\textbf{Non-homogenous} once again exclude a singularity and find a Dirac Delta function.
		
		\textbf{Maximum Principle of Parabolic boundary $\Gamma$}
		
		\textbf{Maximum Principle on an unbounded domain} More complex.  Must occur during initial conditions.  Proof is more complicated (probably not on test).
		
		\textbf{Maximum implies Uniqueness}
		
		\textbf{Regularity estimates:}
		
		Homogenous: have a solution for unbounded case, not bounded case.
		
		Unbounded: Didn't do Green's Function for Hyperbolic Equation (Heat) but you can.  Use cut off function and a double convolustion.
		
		Regularity for $u$ is defined $\dots$.
		
		\textbf{Energy Methods on bounded domain}  
		\begin{itemize}
			\item multiply by something and integrate by parts.  The most popular is $u$ itself.  This gives you a boundary term which typically goes to zero from boundary conditions.  Expected to stay the same or disipate over time.
		\end{itemize}
	
	\item \textbf{Parabolic PDE: Wave Equation $u_{tt}-\Delta_u=0$}
	
	various solution formulas: $(u_t-u_x)(u_t+u_x)=0$ and notice that this is true with the left most term $u_t-u_x=0$.  Speed limit is $c=1$.
	
	Finite speed of propogation (from the solution formula).
	
	d'Alembert's Formula.
	
	\textbf{Regularity Issues} One dimensional case, the three dimensional case is easier that two dimensions.

	\begin{itemize}
	\item	Three dimensions: method of spherical means (two dimensions sets $z\equiv 0$.  Translates to a one-dimensional wave equation on the radius ($r>0$, gotten by using odd reflection on the entire real line).
	
	\item Three dimensional wave propagation continues past the point.
	
	\item Two dimensional case use method of descent.  End up with a cumbersom formula.
	
\end{itemize}		
	
	\textbf{Huygen's Principle} is the difference between 2-dimensional and 3-dimensional wave propogation.
	
	\textbf{Energy Methods} Prove:
	\begin{itemize}
		\item unique
		\item no maximum principle
		\item alternate proof of finite speed of popagation (doesn't use Kirchoff's formula).
	\end{itemize}
		
	\item \textbf{General transport equation}
	\begin{itemize}
		\item Variable-coefficient linear transport
		
		still use method of characteristics \textbf{\textit{•Know this for this test: various things that could happen with method of characteristics.}} $Z(s)=u(x(s), t(s))$ then take the derivative and choose how $\frac{dx}{ds}$ and $\frac{dt}{ds}$ make sense for a solution.
		
		\item non-linear transport, i.e., $u_t-u_x=u^2$
		
		could be a question understanding definition between linear and non-linear PDE.  Beware linear 'operator' (like the Laplacian).
		

\textbf{Strategy: Method of Characteristics}

The method of characteristics transforms the PDE into a system of ordinary differential equations (ODEs) along specific curves in the \((x,t)\) plane, called characteristic curves. Along these curves, the solution's behavior simplifies dramatically.
\begin{enumerate}

\item  Set up the Characteristic Equations

For a general first-order PDE of the form \(a(x,t,u)u_x + b(x,t,u)u_t = c(x,t,u)\), the characteristic equations are:

\[ \frac{dx}{ds} = a, \quad \frac{dt}{ds} = b, \quad \frac{du}{ds} = c \]

For your equation, \(u_t - u_x = u^2\), we can write it as \((-1)u_x + (1)u_t = u^2\). So, \(a=-1\), \(b=1\), and \(c=u^2\). The characteristic equations are:

\begin{enumerate}
\item  \(\frac{dt}{ds} = 1\)
\item  \(\frac{dx}{ds} = -1\)
\item \(\frac{du}{ds} = u^2\)

\end{enumerate}

\item  Solve the ODEs

We solve this system of ODEs, typically with initial conditions given along a curve, for instance, \(u(x,0) = g(x)\) at \(t=0\).

*   **Solve for t:** From \(\frac{dt}{ds} = 1\), we get \(t(s) = s + C\). If we start the characteristics at \(t=0\), we can set \(t=s\).

*   **Solve for x:** From \(\frac{dx}{ds} = -1\), we get \(x(s) = -s + x_0\), where \(x_0\) is the initial x-position at \(s=0\) (i.e., \(t=0\)). Substituting \(s=t\), we find the equation for the characteristic curves:
    \[ x = -t + x_0 \implies x_0 = x+t \]
    This tells us that the characteristic curves are straight lines along which the quantity \(x+t\) is constant.

*   **Solve for u:** Now we solve the crucial ODE for \(u\):
    \[ \frac{du}{ds} = u^2 \]
    This is a separable ODE: \(\frac{du}{u^2} = ds\). Integrating gives:
    \[ -\frac{1}{u} = s + K \implies u(s) = -\frac{1}{s+K} \]
    The constant of integration \(K\) depends on the specific characteristic curve, which is identified by its starting point \(x_0\). At \(s=0\) (which corresponds to \(t=0\)), the initial value of \(u\) is \(u(0) = g(x_0)\). So:
    \[ u(0) = g(x_0) = -\frac{1}{K} \implies K = -\frac{1}{g(x_0)} \]

\item Assemble the Final Solution

Now, substitute everything back to express the solution \(u\) in terms of \(x\) and \(t\).

\[ u(x,t) = u(s) = -\frac{1}{s - \frac{1}{g(x_0)}} \]

Replace \(s\) with \(t\) and \(x_0\) with \(x+t\):

\[ u(x,t) = -\frac{1}{t - \frac{1}{g(x+t)}} = \frac{g(x+t)}{1 - t \cdot g(x+t)} \]

This is the general solution for the initial value problem. A key feature of this non-linear equation is the possibility of the solution blowing up in finite time when the denominator becomes zero, i.e., when \(t \cdot g(x+t) = 1\). This is a stark contrast to the linear PDEs we discussed earlier, whose solutions typically exist for all time (given smooth data).
\end{enumerate}	
		
	\end{itemize}
	
	\item \textbf{Burger's Equation, Conservation laws.}  $u_t+(F(u))_x=u_t-F'(u)u_x=0$.  \textbf{Likely to be on the final}
		
		Start the same way with $z(s)$ and method of characteristics.
		
		could have rarefaction wave and a shock (probably only one at a time).  jump conditions.  Non-uniqueness issue because the characteristics cross.  Return to the weak formulation of the PDE.  \textbf{probably not on the final.}  
		
		\begin{description}
			\item \textbf{Entropy condition.} violating is a non-physical shock.  Since the Method of Characteristics can generate a family of solutions and that these can travel at different speeds (some catching up and overpowering others) they can form a shock.    In this case, many solutions aren't practical or cannot represent a physical solution without the entropy condition.  If $F$ is the characteristic speed of a weak-solution then
		\begin{align*}
			F'(u_L)> s > F'(u_R)
		\end{align*}
		\end{description}
	
	\item \textbf{The solution space, eigenvalues and eigen-functions}
	
	Fundamentally, think of the PDE as $\Delta u = \lambda u$ and to think of the solution space as an inner product space over continuous functions and the inner product is 
	\begin{align*}
		f \cdot g = \int_{\R^n} fg\, dx
	\end{align*}thus we can approach the solution as a linear algebra problem where $\lambda$ is an eigenvalue for a transformation equivalent to the PDE.
	
	\begin{enumerate}
		\item Assume that it is separable $u(x,t) = T(t)\phi(x)$.  Then, in hyperbolic PDEs we have 
		\begin{align*}
			u(x,t) &= T(t)\phi(x) \\
			T'(t)\phi(x) &= T(t)\Delta \phi \\
			\frac{T'(t)}{T(t)} &= \frac{\Delta \phi(x)}{\phi (x)}
		\end{align*}This must be a constant which we assign to $\lambda$.  Generating two equations
		\begin{align*}
			\Delta \phi &= -\lambda \phi & \text{spatial part}\\
			&\text{generates } \{\lambda_n\} \\
			T'(t) &= -\lambda T(t) & \text{time part} \\
			&= C_ne^{\lambda_nt}
		\end{align*}Thus, applying $t=0$ we get
		\begin{align*}
			f(x) &= u(x,0) = \sum_{n=1}^\infty C_n\phi_n(x) \\
			C_n &= \frac{f\cdot\phi}{\phi\cdot\phi}= \frac{\int_U f(x)\phi(x)\,dx}{\int_U \phi_n(x)^2\, dx}
		\end{align*}generating a Green's Function
		\begin{align*}
			G(x,t) &= \sum_{n=1}^\infty \frac{\phi_n(x)\phi_n(t)}{\lambda_n}
		\end{align*}
	\end{enumerate}
	
	\end{itemize}

\end{description}

\end{document}
